\documentclass[conference]{IEEEtran}

\usepackage[utf8]{inputenc}
\usepackage{amsmath}
\usepackage{graphicx}
\usepackage{booktabs}
\usepackage{multirow}
\usepackage{url}
\usepackage{amsfonts}
\usepackage{cite}
\usepackage{xcolor}
\usepackage{float}

% correct bad hyphenation here
\hyphenation{op-tical net-works semi-conduc-tor}

\begin{document}

\title{Art Style Classification Using Hybrid CNN with Edge Detection: A Comparative Study}

\author{
\IEEEauthorblockN{Selman Turan Toker}
\IEEEauthorblockA{AI and Data Engineering\\
Istanbul Technical University\\
Istanbul, Turkey\\
150220330 \\
toker22@itu.edu.tr}
\and
\IEEEauthorblockN{M. Burak KORKMAZ}
\IEEEauthorblockA{AI and Data Engineering\\
Istanbul Technical University\\
Istanbul, Turkey\\
150220326 \\
korkmazmu22@itu.edu.tr}
}

\maketitle

\begin{abstract}
Classifying art styles automatically is a challenging task in computer vision because the differences between styles can be subtle, while paintings within the same style can exhibit significant variation. In this paper, we propose a novel hybrid Convolutional Neural Network (CNN) architecture designed to classify art styles more accurately by leveraging both RGB color information and structural edge features. We evaluated our model on the WikiArt dataset using 20 distinct art styles, ranging from the Renaissance to modern movements. Our architecture integrates a pretrained ResNet50 for RGB feature extraction and a custom CNN stream that processes Canny edge maps to capture brushwork and composition. We compared our approach against two strong baselines: VGG16 and Vision Transformer (ViT). Experimental results demonstrate that the hybrid model achieves approximately 60.2\% accuracy, outperforming the VGG16 baseline 56.2\% and remaining highly competitive with ViT 60.5\%. Furthermore, Grad-CAM visualizations confirm that incorporating edge information helps the model focus on structural boundaries, improving interpretability. Our findings suggest that structural features are essential for robust art style recognition.
\end{abstract}

\IEEEpeerreviewmaketitle

\section{Introduction}

Art style classification is a complex problem that bridges the gap between computer vision and art history. Successfully automating this task has practical applications in digital archiving, fraud detection, and recommendation systems for museums. Unlike general object classification, recognizing art styles requires understanding abstract concepts such as brushwork techniques, color palettes, and structural composition.

Several factors make this task particularly difficult: (1) \textbf{Subtle Differences}: Styles originating from the same era often share similar visual characteristics; (2) \textbf{High Intra-class Variation}: A single movement can include vastly different artistic approaches; (3) \textbf{Data Scarcity}: Annotated art datasets are significantly smaller than standard datasets like ImageNet; and (4) \textbf{Domain Shift}: Models trained on high-quality scans often struggle with real-world photographs taken in museums.

While traditional methods relied on handcrafted features, deep learning has become the standard approach. However, standard CNNs often treat artworks as generic images, focusing heavily on texture and color. Standard RGB models often overlook the structural details of lines and brushstrokes, yet these edge features are fundamental to distinguishing art styles.

In this project, we propose a \textbf{hybrid CNN architecture} that processes two distinct inputs: (1) standard RGB images analyzed by ResNet50 to capture color and texture, and (2) Canny edge maps processed by a custom CNN to extract structural features. These two streams are fused to produce the final classification.

\subsection{Research Questions}

This study aims to answer the following questions:
\begin{enumerate}
    \item Can the explicit inclusion of edge detection features improve classification accuracy compared to RGB-only models?
    \item How does the proposed hybrid architecture compare to state-of-the-art models like Vision Transformers?
    \item Which specific art styles benefit the most from structural edge information?
    \item How robust is the model when tested on cross-domain data collected from real-world museum environments?
\end{enumerate}

\section{Related Work}

\subsection{Traditional Art Analysis}
Before the deep learning era, researchers relied on manual feature extraction. Techniques such as color histograms and Local Binary Patterns (LBP) were used to analyze brushstrokes and texture. While these methods provided some insights, they generally failed to capture the complexity of modern art movements.

\subsection{Deep Learning for Art}
With the rise of CNNs, Transfer Learning became a popular technique. Models pretrained on ImageNet, such as VGG and ResNet, demonstrated significantly better performance than traditional methods. Recent studies have also explored graph neural networks and attention mechanisms to focus on stylistically relevant regions of an artwork.

\subsection{Vision Transformers}
Vision Transformers (ViT) have recently set new benchmarks in image classification by using self-attention mechanisms to analyze global context. Although ViT has shown promise in art classification, these models are typically computationally expensive and require large amounts of data to train effectively.

\subsection{Multi-Stream Networks}
Multi-stream architectures are widely used in video analysis. However, the use of explicit edge detection as a secondary stream for art classification remains relatively unexplored. Our project seeks to address this gap.

\section{Dataset}

\subsection{WikiArt Dataset}
We utilized the WikiArt dataset, sourced from Kaggle, as our primary training source. To ensure a manageable scope, we selected 20 representative art styles spanning various historical periods:

\textbf{Renaissance/Early Modern:} Early Renaissance, High Renaissance, Mannerism, Northern Renaissance, Baroque, Rococo.

\textbf{19th Century:} Romanticism, Realism, Impressionism, Post-Impressionism, Symbolism.

\textbf{Modern/Contemporary:} Art Nouveau, Fauvism, Expressionism, Cubism, Abstract Expressionism, Naive Art, Pop Art, Minimalism, Color Field Painting.

\subsection{Preprocessing}
After filtering the dataset, we obtained approximately \textbf{78,545 images}. To balance the dataset and reduce computational load, we capped the training set at 1,500 images per class. The data was split into 80\% for training, 10\% for validation, and 10\% for testing. All images were resized to $224 \times 224$ pixels. We also applied data augmentation techniques, including rotation, cropping, and horizontal flipping, to prevent overfitting.

\subsection{Museum Dataset (Cross-Domain)}
To evaluate the model's performance in real-world scenarios, we curated a separate \textbf{Museum Dataset}. This dataset consists of photos taken from various museums and exhibitions. Unlike the clean, scanned images in WikiArt, these photographs present significant challenges. Factors such as varying lighting conditions, shadows cast by frames, glare on the canvas, and perspective distortion due to camera angles create a difficult testing environment.

% FIGURE 1: Dataset Examples
\begin{figure}[htbp]
\centering
\includegraphics[width=0.9\columnwidth]{dataset_examples.png}
\caption{Representative samples from the WikiArt dataset showing diverse styles.}
\label{fig:dataset}
\end{figure}

\section{Methodology}

\subsection{Architecture Overview}
Our proposed model is a dual-stream network comprising an RGB stream, an Edge stream, and a fusion block. The overall structure is illustrated in Figure \ref{fig:architecture}.

% FIGURE 2: Architecture
\begin{figure*}[htbp]
\centering
\includegraphics[width=0.8\textwidth]{hybrid_architecture.png}
\caption{The proposed Hybrid CNN architecture. It processes RGB images and Edge maps in parallel streams before fusing them for classification.}
\label{fig:architecture}
\end{figure*}

\subsection{RGB Stream: ResNet50}
For the RGB input stream, we employ a ResNet50 model pretrained on ImageNet. Since pretrained weights are already capable of recognizing complex shapes and textures, we utilize the feature map from the final convolutional layer. This results in a 2048-dimensional feature vector that encapsulates high-level semantic information.

\subsection{Edge Detection Stream}
This stream focuses specifically on the structural components of the artwork. First, we convert the input images to grayscale and apply the Canny edge detection algorithm with thresholds set to 100 and 200. Since edge maps are visually simpler than RGB images, we designed a lightweight custom CNN for this stream. The edge stream network consists of three convolutional layers with 64, 128, and 256 filters respectively. A Max Pooling layer is applied after the second convolutional layer, and an Adaptive Average Pooling layer is used at the end of the stream.

% FIGURE 3: Edge Examples
\begin{figure}[htbp]
\centering
\includegraphics[width=0.9\columnwidth]{edge_examples.png}
\caption{Visualizing the Canny edge detection inputs. The edge maps highlight brushstrokes and geometric forms.}
\label{fig:edges}
\end{figure}

\subsection{Fusion and Classification}
The feature vectors from both streams are concatenated to form a combined vector of size 2304. This fused vector is then passed through a series of fully connected dense layers. We apply ReLU activation functions and Dropout layers to mitigate overfitting. Finally, a Softmax layer is used to predict the probability distribution across the 20 art style classes.

\subsection{Baseline Models}
To provide a fair evaluation of our hybrid approach, we compared it against two distinct baselines:
\begin{itemize}
    \item \textbf{VGG16:} A classic, standard CNN architecture.
    \item \textbf{Vision Transformer (ViT):} A modern attention-based architecture.
\end{itemize}

\subsection{Training Details}
We trained models using the Adam optimizer. For VGG16 and ViT, we used a learning rate of 0.0001. For the Hybrid CNN, we used differential learning rates: 1e-5 for the RGB stream, 8e-4 for the edge stream, and 5e-4 for the classifier. We utilized CrossEntropy loss with Label Smoothing set to 0.15. The training was conducted with a batch size of 64 for 30 epochs, employing Early Stopping to prevent overfitting.

\section{Experiments and Results}

\subsection{Metrics}
We evaluated the models using Top-1 Accuracy, Top-5 Accuracy, and the F1-Score (Macro average). The F1-Score is particularly important in this context to ensure that performance is not biased toward classes with slightly more data.

\subsection{Quantitative Results}
First, we evaluated the performance of the proposed Hybrid model against the baselines on the standard WikiArt test set. The results are summarized in Table \ref{tab:test_performance}.

\begin{table}[H]
\centering
\caption{Test Set Performance Comparison (WikiArt)}
\label{tab:test_performance}
\begin{tabular}{@{}lccc@{}}
\toprule
Model & Top-1 Acc (\%) & Top-5 Acc (\%) & F1 Score (\%) \\ \midrule
Hybrid CNN & \textbf{59.02} & 100.0 & \textbf{58.68} \\
VGG16 & 54.33 & 100.0 & 53.24 \\
ViT (Base) & 58.29 & 100.0 & 58.38 \\ \bottomrule
\end{tabular}
\end{table}

As seen in Table \ref{tab:test_performance}, our Hybrid CNN outperforms the VGG16 baseline by approximately 4.7\%, demonstrating the significant value of adding explicit edge information. Furthermore, it achieves results slightly superior to the Vision Transformer (59.02\% vs 58.29\%), which is a significantly more complex architecture.

To assess the generalization capability of our approach, we further evaluated the models on a combined dataset that integrates the WikiArt test set with our real-world Museum dataset. Table \ref{tab:dataset_comparison} compares the Top-1 accuracy across these settings.

\begin{table}[H]
\centering
\caption{Model Comparison: WikiArt vs. Combined Dataset}
\label{tab:dataset_comparison}
\begin{tabular}{@{}lcc@{}}
\toprule
Model & WikiArt Top-1 (\%) & Combined Top-1 (\%) \\ \midrule
Hybrid CNN & \textbf{59.0} & \textbf{60.2} \\
VGG16 & 54.3 & 56.2 \\
ViT & 58.3 & 60.5 \\ \bottomrule
\end{tabular}
\end{table}

When tested on the pure Museum dataset (which constitutes the difference in the combined set), all models experienced a performance fluctuation due to domain shifts such as lighting and perspective. However, as shown in Table \ref{tab:dataset_comparison}, the Hybrid CNN maintained its performance lead over the VGG16 baseline in the combined scenario, indicating reasonable robustness against real-world variations.

% FIGURE 4: Training History
\begin{figure}[htbp]
\centering
\includegraphics[width=0.9\columnwidth]{training_history.png}
\caption{Training and validation accuracy curves. The hybrid model demonstrates stable learning progression.}
\label{fig:training}
\end{figure}

\subsection{Confusion Matrix Analysis}
An analysis of the confusion matrix reveals distinct patterns in the model's errors. As expected, the model sometimes confuses stylistically similar movements, such as \textit{Impressionism} and \textit{Post-Impressionism}. However, styles with distinct geometric structures, such as \textit{Cubism} and \textit{Minimalism}, are classified with high accuracy. This strongly suggests that the edge stream is effectively capturing the necessary structural features.

% FIGURE 5: Confusion Matrix
\begin{figure}[htbp]
\centering
\includegraphics[width=0.9\columnwidth]{confusion_matrix.png}
\caption{Confusion matrix for the Hybrid CNN. Darker diagonals indicate correct predictions.}
\label{fig:confusion}
\end{figure}

\subsection{Grad-CAM Visualization}
We employed Grad-CAM to visualize the regions of interest for each model. As shown in Figure \ref{fig:gradcam}, VGG16 primarily focuses on color distributions. In contrast, the Hybrid CNN attends to both color patches and structural lines. This dual focus allows the hybrid model to make more informed and interpretable decisions.

% FIGURE 6: Grad-CAM
\begin{figure}[htbp]
\centering
\includegraphics[width=0.9\columnwidth]{gradcam_comparison.png}
\caption{Grad-CAM visualizations. The Hybrid model (bottom row) focuses on structural boundaries more effectively than VGG16.}
\label{fig:gradcam}
\end{figure}
\section{Conclusion}

In this project, we successfully developed and evaluated a hybrid Deep Learning architecture for art style classification. By fusing standard RGB features from ResNet50 with structural features from a custom edge-detection CNN, we demonstrated that multi-modal input can enhance model performance.

Our experiments yielded several key insights. First, the explicit inclusion of edge information improves accuracy compared to standard CNNs like VGG16. Second, our hybrid model remains competitive with state-of-the-art architectures like Vision Transformers while being potentially more interpretable. Finally, we observed that art movements defined by strong geometric shapes, such as Cubism, benefit significantly from edge detection.

The project highlights that while deep learning models are powerful, incorporating domain specific knowledge can lead to better and more robust systems. Future work could explore learnable edge detection layers to replace fixed Canny thresholds or implement more advanced attention mechanisms to dynamically weight the RGB and Edge streams.

\begin{thebibliography}{00}

\bibitem{b1} S. Karayev et al., ``Recognizing Image Style,'' \textit{BMVC}, 2014.
\bibitem{b2} A. Elgammal et al., ``Picasso, Matisse, or a Fake?,'' \textit{AAAI}, 2018.
\bibitem{b3} K. He et al., ``Deep Residual Learning for Image Recognition,'' \textit{CVPR}, 2016.
\bibitem{b4} K. Simonyan and A. Zisserman, ``Very Deep Convolutional Networks,'' \textit{ICLR}, 2015.
\bibitem{b5} A. Dosovitskiy et al., ``An Image is Worth 16x16 Words: Transformers for Image Recognition at Scale,'' \textit{ICLR}, 2021.
\bibitem{b6} J. Canny, ``A Computational Approach to Edge Detection,'' \textit{IEEE Trans. PAMI}, 1986.
\bibitem{b7} R. R. Selvaraju et al., ``Grad-CAM,'' \textit{ICCV}, 2017.
\bibitem{b8} WikiArt Dataset, Kaggle: \url{https://www.kaggle.com/steubk/wikiart}

\end{thebibliography}

\end{document}